% =====================================================================
\chapter{Introduction}\label{sec:intro}
% =====================================================================

\section{Motivation}
In a normal Internet of Vehicles (IoV), roadside units (RSUs) on poles verify every vehicle and relay data to a cloud Trusted Authority (TA).
After an earthquake or flood the poles, the fibre and the power are the first things to fail. Rescue convoys and evacuation traffic then need secure
coordination most, and nobody is left to verify anyone. Tan, Zheng and Vijayakumar~\cite{base} call this the \emph{infrastructure-less IoV}.
Their answer is a \emph{tethered UAV} (TUAV): a drone powered by cable from an emergency communications response vehicle (ECRV), which hovers for hours
and acts as a flying edge server for a swarm of ordinary UAVs.

\section{The base paper}
The base paper makes three contributions:
\begin{enumerate}[label=\textbf{(\roman*)}]
  \item a TUAV-enabled edge architecture that replaces the dead RSU and gives three-dimensional coverage;
  \item certificateless \emph{group} authentication of $n$ UAVs at once, with the costly bilinear pairings offloaded to the TA;
  \item a group key distributed with the Chinese Remainder Theorem (CRT) and updated cheaply when UAVs join or leave.
\end{enumerate}
It reports 151\,ms to authenticate 120 UAVs, 104 bytes per request, and all eleven security properties of its comparison table.

\section{Problem statement}
\begin{keybox}[title={Research question}]
Does the published scheme actually work \emph{without infrastructure}, and does it deliver the security properties it claims? If not, can those
properties be recovered together, at comparable cost, and can the same design be extended to the vehicles and to a multi-TUAV deployment?
\end{keybox}

\section{Summary of findings}
We re-derived every equation of the base paper, implemented it on a real pairing-friendly curve, and modelled its credential in ProVerif. The
architecture is a good idea, but the protocol has six problems (Chapter~\ref{sec:findings}):
\begin{enumerate}[label=\textbf{F\arabic*.}]
  \item \textbf{It is not infrastructure-less.} The TUAV cannot verify anyone itself; every authentication needs a satellite round trip to the TA.
  \item \textbf{The TA can impersonate every UAV} (no escrow resilience, no non-repudiation). Our code reproduces this, and ProVerif finds it automatically.
  \item \textbf{No forward secrecy}: one leaked long-term secret exposes every past group key.
  \item \textbf{Revocation does not revoke}: the CRT polynomial reveals every member's key, so a removed UAV still decrypts. Our code reproduces this.
  \item \textbf{The unforgeability proof is empty}, and the batch check accepts cancelling errors. Our code reproduces the cancellation.
  \item \textbf{The cost accounting omits the TA's pairings} and an $O(n^3)$-byte key broadcast. We measured both.
\end{enumerate}

\section{Contributions}
\begin{keybox}
Keep the TUAV architecture and move the trust root \emph{offline}. The TA certifies credentials before deployment and can trace a misbehaving entity
afterwards, but no operation in the field ever waits for it. Every UAV, vehicle and TUAV can verify every other one using only the TA's public key.
\end{keybox}
\begin{enumerate}[label=\textbf{\arabic*.},leftmargin=*]
  \item \textbf{Analysis} of the base scheme: six findings, four reproduced in code and one confirmed by ProVerif (Chapter~\ref{sec:findings}).
  \item \textcolor{uavc}{\textbf{C1, autonomous UAV tier}}: escrow-free, pairing-free certificateless credentials verified locally by the TUAV in a
        small-exponent batch with a bounded fallback; per-member group-key delivery under fresh ECDH keys; secure rekeying (Chapter~\ref{sec:c1}).
  \item \textcolor{vehc}{\textbf{C2, vehicle tier}}: vehicles admitted through authenticated UAVs that relay and filter their requests, with a separate
        regional key for V2V (Chapter~\ref{sec:c2}).
  \item \textcolor{hoc}{\textbf{C3, multi-TUAV handover}}: single-use encrypted tickets and one fresh ECDH, with no signature verification and no TA (Chapter~\ref{sec:c3}).
  \item \textbf{Formal verification} in ProVerif: five models and sixteen queries, including the base scheme's escrow attack and the exact boundary of our
        own escrow guarantee (Chapter~\ref{ch:pv}).
  \item \textbf{Implementation and measured evaluation}: both schemes on the same curve and library, 20 tests, and eight experiments covering
        availability, latency, batch poisoning, traffic, vehicles, handover under mobility, and attacks (Chapters~\ref{ch:impl}--\ref{ch:eval}).
\end{enumerate}

\section{Headline results}
\begin{center}\small
\begin{tabularx}{\linewidth}{@{}>{\raggedright\arraybackslash}X cc l@{}}
\toprule
Metric & Base scheme & Ours & Evidence\\
\midrule
Authentication when the TA link is down & 0\% & 100\% & measured model (E1)\\
TA can impersonate a registered UAV & yes & no & ProVerif + code (E7)\\
Revoked UAV obtains the new group key & yes & no & ProVerif + code (E7)\\
Honest UAVs admitted with one bad request & 0\% & all & code (E3, E7)\\
Group-key traffic, $n=120$ & 17.5\,MB & 9.6\,KB & measured (E4)\\
Authenticate + key 120 UAVs, base paper's constants & 3720\,ms + RTT & 228\,ms & analytical\\
\bottomrule
\end{tabularx}
\end{center}

\section{Organisation}
Chapter~\ref{ch:background} reviews the background. Chapter~\ref{sec:base} recalls the base scheme and Chapter~\ref{sec:findings} analyses it.
Chapter~\ref{sec:model} gives the system and threat model, the design goals and the notation. Chapters~\ref{sec:setup}--\ref{sec:tr} present the framework,
Chapter~\ref{sec:security} the security analysis and Chapter~\ref{ch:pv} the formal verification. Chapter~\ref{ch:impl} describes the implementation.
Chapters~\ref{sec:perf} and~\ref{ch:eval} report the analytical and the measured evaluation. Chapter~\ref{ch:conclusion} discusses limitations and concludes.
